\documentclass[12pt]{article}
\usepackage[utf8]{inputenc}
\usepackage[T1]{fontenc}
\usepackage[margin=1in]{geometry}
\usepackage{setspace}
\usepackage{xcolor}
\usepackage[colorlinks=true, urlcolor=blue!60!black, linkcolor=blue!60!black]{hyperref}
\setlength{\parindent}{2em}
\setlength{\parskip}{1.25em}
\renewcommand{\baselinestretch}{1.0}
\begin{document}
\begin{center}
{\LARGE \textit{Lemon v. Kurtzman}, 403 U.S. 602 (1971)\par}
\vspace{0.8em}
{\large Case Brief\par}
\end{center}
\vspace{1em}

\section{Memory Jogger}
The Supreme Court struck down Pennsylvania and Rhode Island programs funding secular instruction in parochial schools because policing those subsidies created excessive government entanglement with religion, synthesizing the famous three-prong \textit{Lemon} test (secular purpose, primary effect, and no excessive entanglement) that governed Establishment Clause doctrine for decades before being abandoned in 2022.

\section{Detailed Case Facts}
\textit{Lemon v. Kurtzman} consolidated challenges to two state statutes enacted to alleviate financial pressures on nonpublic elementary and secondary schools, the vast majority of which were Roman Catholic parochial schools (\href{https://www.courtlistener.com/opinion/108380/lemon-v-kurtzman/}{CourtListener}; \href{https://supreme.justia.com/cases/federal/us/403/602/}{Justia}).

\textbf{Rhode Island's 1969 Salary Supplement Act} (R.I. Gen. Laws Ann. \S~16-51-1 et seq.) authorized state officials to pay a salary supplement of up to 15\% directly to teachers of secular subjects in nonpublic elementary schools where average per-pupil expenditures on secular education fell below the public-school average (\href{https://supreme.justia.com/cases/federal/us/403/602/}{403 U.S. at 607--08}). Eligible teachers had to be state-certified, teach only subjects offered in public schools using only public-school materials, and agree in writing not to teach any course in religion while receiving the supplement (\href{https://supreme.justia.com/cases/federal/us/403/602/}{403 U.S. at 607--08}). Approximately 25\% of Rhode Island's elementary students attended nonpublic schools, roughly 95\% of whom attended Roman Catholic schools; all 250 teachers who applied for benefits taught in Roman Catholic schools governed by the Bishop of Providence and a diocesan ``Handbook of School Regulations'' (\href{https://supreme.justia.com/cases/federal/us/403/602/}{403 U.S. at 608, 617--18}).

\textbf{Pennsylvania's 1968 Nonpublic Elementary and Secondary Education Act} (Pa. Stat. Ann., Tit. 24, \S\S~5601--5609) authorized the state Superintendent of Public Instruction (David H. Kurtzman) to ``purchase'' secular educational services by directly reimbursing nonpublic schools for teachers' salaries, textbooks, and instructional materials in four secular subjects: mathematics, modern foreign languages, physical science, and physical education (\href{https://supreme.justia.com/cases/federal/us/403/602/}{403 U.S. at 609--10}). Textbooks and materials required state approval, courses containing religious teaching or sectarian morals were barred from reimbursement, and participating schools had to maintain state-audited accounting records separating secular and religious costs (\href{https://supreme.justia.com/cases/federal/us/403/602/}{403 U.S. at 609--10}). Funded initially by horse- and harness-racing taxes and later by cigarette taxes, the program spent about \$5 million annually across 1,181 nonpublic schools serving 535,215 students---more than 20\% of Pennsylvania's students, over 96\% of whom attended church-related (mostly Roman Catholic) schools (\href{https://supreme.justia.com/cases/federal/us/403/602/}{403 U.S. at 610}).

\section{Procedural History}
In the Rhode Island litigation (\textit{Earley v. DiCenso}, No. 569, and \textit{Robinson v. DiCenso}, No. 570), Rhode Island citizens and taxpayers sued state education officials, eligible teachers, and parochial-school parents in federal court to enjoin the Salary Supplement Act under the First Amendment's Religion Clauses (\href{https://supreme.justia.com/cases/federal/us/403/602/}{403 U.S. at 608}). After an evidentiary hearing, a three-judge U.S. District Court for the District of Rhode Island held the Act unconstitutional under the Establishment Clause because it fostered ``excessive entanglement'' between government and religion, 316 F. Supp. 112 (D.R.I. 1970) (\href{https://supreme.justia.com/cases/federal/us/403/602/}{403 U.S. at 608--09}).

In the Pennsylvania litigation (\textit{Lemon v. Kurtzman}, No. 89), individual citizens and taxpayers---including Alton J. Lemon, a public-school parent who alleged he had paid the racetrack tax supporting the Act---alongside separationist organizations sued state officials and seven church-related schools (\href{https://supreme.justia.com/cases/federal/us/403/602/}{403 U.S. at 610--11}). A three-judge U.S. District Court for the Eastern District of Pennsylvania recognized individual taxpayer standing under \textit{Flast v. Cohen}, denied organizational standing, and granted the defendants' motion to dismiss for failure to state a claim, holding over Chief Judge Hastie's dissent that the Act violated neither the Establishment nor the Free Exercise Clause, 310 F. Supp. 35 (E.D. Pa. 1969) (\href{https://supreme.justia.com/cases/federal/us/403/602/}{403 U.S. at 611}). On direct appeal, the U.S. Supreme Court consolidated Nos. 89, 569, and 570, heard oral argument on March 3, 1971, and decided the cases on June 28, 1971 (\href{https://supreme.justia.com/cases/federal/us/403/602/}{403 U.S. at 602, 606}; \href{https://www.oyez.org/cases/1970/89}{Oyez}). (On remand in Pennsylvania, disputes over reimbursement for services already rendered before the 1971 decision produced \href{https://supreme.justia.com/cases/federal/us/411/192/}{\textit{Lemon v. Kurtzman} (\textit{Lemon II}), 411 U.S. 192 (1973)}.)

\section{Judicial Votes}
\begin{itemize}
\item \textbf{Vote in No. 89 (Pennsylvania statute):} 8--0, reversing and remanding (Justice Thurgood Marshall took no part in the consideration or decision of No. 89) (\href{https://supreme.justia.com/cases/federal/us/403/602/}{403 U.S. at 604, 625}).
\item \textbf{Vote in Nos. 569 \& 570 (Rhode Island statute):} 8--1, affirming (\href{https://supreme.justia.com/cases/federal/us/403/602/}{403 U.S. at 604, 625}).
\item \textbf{Majority Opinion:} Chief Justice Warren E. Burger (author), joined by Justices Hugo L. Black, William O. Douglas, John Marshall Harlan II, Potter Stewart, Harry A. Blackmun, and (as to Nos. 569 and 570 only) Thurgood Marshall (\href{https://supreme.justia.com/cases/federal/us/403/602/}{403 U.S. at 604}).
\item \textbf{Concurrence (Douglas):} Justice William O. Douglas (author), joined by Justice Hugo L. Black, and joined as to Nos. 569 and 570 by Justice Thurgood Marshall (who added a separate statement reserving judgment on \textit{Everson v. Board of Education}) (\href{https://supreme.justia.com/cases/federal/us/403/602/}{403 U.S. at 604, 625, 642}).
\item \textbf{Concurrence (Brennan):} Justice William J. Brennan, Jr. (author), concurring in Nos. 89, 569, and 570 (and dissenting in companion case \textit{Tilton v. Richardson}, 403 U.S. 672 (1971)) (\href{https://supreme.justia.com/cases/federal/us/403/602/}{403 U.S. at 604, 642}).
\item \textbf{Concurrence in the Judgment in Part and Dissent in Part (White):} Justice Byron R. White (author), concurring in the judgment in No. 89 (Pennsylvania) and dissenting in Nos. 569 and 570 (Rhode Island) (\href{https://supreme.justia.com/cases/federal/us/403/602/}{403 U.S. at 604, 661}).
\end{itemize}

\section{Holding}
\textbf{Judgment:} Affirmed in Nos. 569 and 570 (invalidating Rhode Island's Salary Supplement Act); reversed and remanded in No. 89 (invalidating Pennsylvania's Nonpublic Elementary and Secondary Education Act) (\href{https://supreme.justia.com/cases/federal/us/403/602/}{403 U.S. at 604, 625}).

\textbf{Rule Announced:} To comply with the Establishment Clause of the First Amendment (applicable to the States through the Fourteenth Amendment), a statute must satisfy three cumulative criteria: (1) it must have a secular legislative purpose; (2) its principal or primary effect must be one that neither advances nor inhibits religion; and (3) it must not foster ``an excessive government entanglement with religion'' (\href{https://supreme.justia.com/cases/federal/us/403/602/}{403 U.S. at 612--13}). To evaluate excessive entanglement, courts examine the character and purposes of the benefited institutions, the nature of the state aid, and the resulting relationship between government and religious authority (\href{https://supreme.justia.com/cases/federal/us/403/602/}{403 U.S. at 615}).

\textbf{Holding vs. Dicta:} The actual holding rested exclusively on the third prong---excessive administrative and political entanglement---because Chief Justice Burger accepted the legislatures' stated secular educational purposes and explicitly declined to decide whether the statutes failed the second (``primary effect'') prong (\href{https://supreme.justia.com/cases/federal/us/403/602/}{403 U.S. at 613--14}). The broader three-part formulation synthesizing prior cases (\textit{Board of Education v. Allen} and \textit{Walz v. Tax Commission}) functioned as the doctrinal framework for half a century, though later decisions eroded its entanglement and effects prongs (\href{https://supreme.justia.com/cases/federal/us/536/639/}{\textit{Zelman v. Simmons-Harris}, 536 U.S. 639 (2002)}) and the Supreme Court formally declared in \href{https://supreme.justia.com/cases/federal/us/597/21-418/}{\textit{Kennedy v. Bremerton School District}, 597 U.S. 507 (2022)}, that it had ``long ago abandoned \textit{Lemon}'' in favor of ``historical practices and understandings.''

\section{Analysis of Opinions}

\subsection{Majority Opinion (Chief Justice Burger)}
Chief Justice Burger candidly observed that the language of the Religion Clauses is ``at best, opaque,'' prohibiting not merely a state church but any law ``respecting'' an establishment of religion---specifically guarding against ``sponsorship, financial support, and active involvement of the sovereign in religious activity'' (\href{https://supreme.justia.com/cases/federal/us/403/602/}{403 U.S. at 612}). Distilling prior precedents, he announced the three-prong \textit{Lemon} test: secular legislative purpose, primary effect neither advancing nor inhibiting religion, and no excessive government entanglement with religion (\href{https://supreme.justia.com/cases/federal/us/403/602/}{403 U.S. at 612--13}).

Accepting both legislatures' stated secular purpose of improving secular education and bypassing the primary-effect prong, Burger invalidated both statutes on excessive entanglement (\href{https://supreme.justia.com/cases/federal/us/403/602/}{403 U.S. at 613--14}). Distinguishing neutral, one-time secular benefits like bus rides (\textit{Everson}) and secular textbooks (\textit{Allen}), he stressed that unlike a textbook whose content can be inspected once, a teacher in a pervasively religious elementary school operates under ecclesiastical authority and cannot easily compartmentalize faith from secular instruction (\href{https://supreme.justia.com/cases/federal/us/403/602/}{403 U.S. at 616--19}). Ensuring compliance with the statutory bans on religious teaching would therefore require ``comprehensive, discriminating, and continuing state surveillance,'' compounded by state audits inspecting school financial records to separate secular from religious expenditures (\href{https://supreme.justia.com/cases/federal/us/403/602/}{403 U.S. at 619--22}). Pennsylvania's program carried the added vice of direct cash payments to church-related schools rather than aid to parents or students (\href{https://supreme.justia.com/cases/federal/us/403/602/}{403 U.S. at 621--22}). Finally, Burger warned that recurring annual legislative appropriations benefiting primarily one denomination would foster ``political division along religious lines''---a core evil the First Amendment was designed to prevent (\href{https://supreme.justia.com/cases/federal/us/403/602/}{403 U.S. at 622--24}).

\subsection{Concurring Opinion (Justice Douglas, joined by Justice Black and in part by Justice Marshall)}
Justice Douglas joined the Court's opinion but wrote separately to advance a stricter separationist rationale rooted in James Madison's \textit{Memorial and Remonstrance} and \textit{Everson v. Board of Education} (\href{https://supreme.justia.com/cases/federal/us/403/602/}{403 U.S. at 625, 633--34, 640}). Tracing the history of Protestant--Catholic school conflicts and detailing the Providence diocesan handbook, Douglas argued that parochial schools are ``an organism living on one budget'' whose \textit{raison d'\^{e}tre} is propagating religious faith (\href{https://supreme.justia.com/cases/federal/us/403/602/}{403 U.S. at 628, 637--41}). Subsidizing secular teachers frees church funds for religious indoctrination through ``sophisticated bookkeeping,'' while policing the grants puts a ``public investigator into every classroom,'' simultaneously threatening both the Establishment and Free Exercise Clauses (\href{https://supreme.justia.com/cases/federal/us/403/602/}{403 U.S. at 627--28, 634, 641}). Justice Marshall joined Douglas's concurrence as to Rhode Island (Nos. 569 and 570) while reserving any view on the continuing vitality of \textit{Everson} (\href{https://supreme.justia.com/cases/federal/us/403/602/}{403 U.S. at 642}).

\subsection{Concurring Opinion (Justice Brennan)}
Justice Brennan concurred in invalidating both state laws (and would have reversed Pennsylvania outright rather than remanding) while dissenting in the companion higher-education case, \textit{Tilton v. Richardson} (\href{https://supreme.justia.com/cases/federal/us/403/602/}{403 U.S. at 642}). Surveying nineteenth-century state constitutional bans on public aid to sectarian schools, Brennan contended that direct tax subsidies to religious schools are qualitatively different from passive tax exemptions (\textit{Walz}) or student-directed benefits (\textit{Everson} and \textit{Allen}) (\href{https://supreme.justia.com/cases/federal/us/403/602/}{403 U.S. at 645--49, 652--57}). He warned that state policing of classrooms induces religious self-censorship and threatens the ``secularization of a creed,'' and concluded that all three statutes impermissibly used ``essentially religious means to serve governmental ends, where secular means would suffice'' (\href{https://supreme.justia.com/cases/federal/us/403/602/}{403 U.S. at 649--51, 658--59}).

\subsection{Concurring in the Judgment in Part and Dissenting in Part (Justice White)}
Justice White dissented in the Rhode Island cases and concurred only in the judgment reversing the motion to dismiss in Pennsylvania (favoring remand for trial rather than facial invalidation) (\href{https://supreme.justia.com/cases/federal/us/403/602/}{403 U.S. at 661, 665, 670--71}). White argued that states may constitutionally finance the separable secular educational function performed by parochial schools even if religion incidentally benefits, noting that the District Court in Rhode Island had expressly found no evidence that participating teachers injected religion into secular classes (\href{https://supreme.justia.com/cases/federal/us/403/602/}{403 U.S. at 663--64, 666--67}). He famously criticized the majority for creating an ``insoluble paradox'': forbidding states from funding secular instruction without anti-religion safeguards, yet striking down those very safeguards as unconstitutional ``entanglement'' (\href{https://supreme.justia.com/cases/federal/us/403/602/}{403 U.S. at 668}). He also faulted the Court for failing to reconcile its rejection of the state school programs with its approval of federal construction grants to religious colleges in \textit{Tilton} (\href{https://supreme.justia.com/cases/federal/us/403/602/}{403 U.S. at 668--69, 671}).

\section{Examples: Future Applications}
\begin{enumerate}
\item \textbf{Same-Side (Unconstitutional under \textit{Lemon}):} A state enacts a ``Parochial STEM Salary Act'' paying 20\% of the salaries of science and math teachers employed by private religious high schools, conditioned on annual classroom observations by state inspectors and state audits of school budgets to verify that no faith-based commentary occurs in subsidized classes. \textit{Why:} Like the Pennsylvania and Rhode Island laws in \textit{Lemon}, direct funding of religious-school teachers combined with ongoing state classroom surveillance and financial auditing creates excessive administrative entanglement.
\item \textbf{Same-Side (Unconstitutional under \textit{Lemon}):} A state legislature passes a statute requiring every public school classroom to display a state-selected religious creed, where the legislative record shows the sole motivation was to promote a particular religious faith. \textit{Why:} Even without entanglement, the law fails \textit{Lemon}'s first prong (no secular legislative purpose) and second prong (primary effect of advancing religion).
\item \textbf{Opposite-Side (Constitutional):} A state creates a neutral scholarship voucher program granting low-income families \$6,000 per child to spend at any participating public, secular private, or religious private school of the parents' choosing, with no state monitoring of classroom religious content. \textit{Why:} As the Court held in \href{https://supreme.justia.com/cases/federal/us/536/639/}{\textit{Zelman v. Simmons-Harris}, 536 U.S. 639 (2002)}, neutral indirect aid directed to religious schools solely through genuine, independent private choice avoids both direct institutional subsidy and the ongoing classroom surveillance that doomed the statutes in \textit{Lemon}.
\item \textbf{Opposite-Side (Constitutional):} A municipality provides fire safety inspections, building-code enforcement, and secular bus transportation to all elementary students attending accredited public and private schools alike. \textit{Why:} Chief Justice Burger expressly recognized in \textit{Lemon} that neutral general services and routine safety inspections are necessary, non-ideological contacts permitted under \textit{Everson} and \textit{Walz} (\href{https://supreme.justia.com/cases/federal/us/403/602/}{403 U.S. at 614, 616}).
\item \textbf{Fence-Sitter (Close Case):} A state sends publicly employed, state-supervised remedial reading specialists onto parochial school campuses to provide federally funded tutoring to low-income students in dedicated rooms stripped of religious symbols, with periodic unannounced visits by public-school supervisors. \textit{Why:} Under \textit{Lemon}'s original presumption that any public employee working on parochial-school grounds risks inculcating religion and requires entangling surveillance, such programs were initially struck down, yet the Court later reversed course in \textit{Agostini v. Felton}, 521 U.S. 203 (1997), rejecting the presumption of pervasive indoctrination and folding entanglement into the effects inquiry (\href{https://supreme.justia.com/cases/federal/us/536/639/}{\textit{Zelman}, 536 U.S. at 649}).
\end{enumerate}

\section{Critique}
\textbf{Originalist and Textualist Critiques:} Originalist jurists and scholars long assailed the \textit{Lemon} test as ahistorical, manipulable, and divorced from the Establishment Clause's original public meaning (\href{https://journals.law.harvard.edu/jlpp/a-little-noted-puzzle-in-religion-law-post-bremerton-gabrielle-girgis/}{Harvard Journal of Law \& Public Policy}; \href{https://texaslawreview.org/applying-history-as-law-the-role-of-historical-facts-in-implementing-constitutional-doctrine/}{Texas Law Review}). Justice Scalia famously mocked \textit{Lemon} in \href{https://supreme.justia.com/cases/federal/us/508/384/}{\textit{Lamb's Chapel v. Center Moriches Union Free School Dist.}, 508 U.S. 384, 398 (1993)}, as ``some ghoul in a late-night horror movie that repeatedly sits up in its grave and shuffles abroad, after being repeatedly killed and buried.'' Echoing Justice White's ``insoluble paradox,'' critics also argued that \textit{Lemon}'s hostility to religious schools penalized religious families exercising their constitutional right to choose parochial education and invited arbitrary invalidation of longstanding public religious traditions---culminating in \href{https://supreme.justia.com/cases/federal/us/597/21-418/}{\textit{Kennedy v. Bremerton School District}, 597 U.S. 507 (2022)}, which declared \textit{Lemon} abandoned in favor of ``historical practices and understandings.''

\textbf{Progressive and Doctrinal Perspectives:} Separationist and progressive scholars defend \textit{Lemon}'s core insight that preventing government sponsorship, financial support, and active entanglement protects both religious minorities and the autonomy of religious institutions themselves (\href{https://houstonlawreview.org/article/92811-unprincipled}{Houston Law Review}; \href{https://sclawreview.org/article/establishment-clause-jurisprudence-and-the-constitutional-limits-on-religion-in-public-schools/}{South Carolina Law Review}). Critics of \textit{Kennedy v. Bremerton}'s abandonment of \textit{Lemon} argue that replacing a structured three-prong standard with an undefined ``history and tradition'' inquiry leaves lower courts with discordant historical sources, elevates Free Exercise and Free Speech claims above nonestablishment, and risks privileging majority faiths at the expense of religious minorities and nonbelievers (\href{https://houstonlawreview.org/article/92811-unprincipled}{Houston Law Review}).

\textbf{Competing Underlying Values:} \textit{Lemon} embodies the enduring constitutional tension between strict church--state separation (preventing taxpayer funding of religious indoctrination, state intrusion into religious institutions, and sectarian political division) and religious neutrality and accommodation (allowing states to support the secular education of all children without discriminating against families who choose religious schools).

\section{Key Quotations}
\begin{quote}
``Every analysis in this area must begin with consideration of the cumulative criteria developed by the Court over many years. Three such tests may be gleaned from our cases. First, the statute must have a secular legislative purpose; second, its principal or primary effect must be one that neither advances nor inhibits religion, \textit{Board of Education v. Allen}, 392 U.S. 236, 243 (1968); finally, the statute must not foster `an excessive government entanglement with religion.' \textit{Walz}, \textit{supra}, at 674.''
--- Chief Justice Warren E. Burger, Opinion of the Court (\href{https://supreme.justia.com/cases/federal/us/403/602/}{403 U.S. at 612--13}; \href{https://www.courtlistener.com/opinion/108380/lemon-v-kurtzman/}{CourtListener}).
\end{quote}

\begin{quote}
``Judicial caveats against entanglement must recognize that the line of separation, far from being a `wall,' is a blurred, indistinct, and variable barrier depending on all the circumstances of a particular relationship.''
--- Chief Justice Warren E. Burger, Opinion of the Court (\href{https://supreme.justia.com/cases/federal/us/403/602/}{403 U.S. at 614}).
\end{quote}

\begin{quote}
``A comprehensive, discriminating, and continuing state surveillance will inevitably be required to ensure that these restrictions are obeyed and the First Amendment otherwise respected. Unlike a book, a teacher cannot be inspected once so as to determine the extent and intent of his or her personal beliefs and subjective acceptance of the limitations imposed by the First Amendment.''
--- Chief Justice Warren E. Burger, Opinion of the Court (\href{https://supreme.justia.com/cases/federal/us/403/602/}{403 U.S. at 619}).
\end{quote}

\begin{quote}
``The surveillance or supervision of the States needed to police grants involved in these three cases, if performed, puts a public investigator into every classroom and entails a pervasive monitoring of these church agencies by the secular authorities. Yet if that surveillance or supervision does not occur, the zeal of religious proselytizers promises to carry the day and make a shambles of the Establishment Clause.''
--- Justice William O. Douglas, Concurring Opinion (\href{https://supreme.justia.com/cases/federal/us/403/602/}{403 U.S. at 627}).
\end{quote}

\begin{quote}
``The Court thus creates an insoluble paradox for the State and the parochial schools. The State cannot finance secular instruction if it permits religion to be taught in the same classroom; but if it exacts a promise that religion not be so taught---a promise the school and its teachers are quite willing and, on this record, able, to give---and enforces it, it is then entangled in the `no entanglement' aspect of the Court's Establishment Clause jurisprudence.''
--- Justice Byron R. White, Concurring in the Judgment in No. 89 and Dissenting in Nos. 569 and 570 (\href{https://supreme.justia.com/cases/federal/us/403/602/}{403 U.S. at 668}).
\end{quote}

\end{document}
